\section{Computational Methodology}
\label{sec:computation}

Tree enumeration uses the Beyer--Hedetniemi canonical level-sequence
algorithm~\cite{BeyerH80}, generating all non-isomorphic trees on $n$ vertices in
$O(n)$ amortized time per tree. The \texttt{synthesizer} engine outputs
trees in graph6 format on \texttt{stdout}, which is piped to the
\texttt{leontovich\_fast} filter. The filter decodes graph6 strings,
precomputes $A^k \mathbf{1}$ for $k = 0, \ldots, n_{\max}$, and evaluates
$\Hom(E_n^{(d)}, H)$ vs.\ $\Hom(P_n, H)$ via matrix--vector iteration in
IEEE 754 double precision. Because the adjacency matrices and state vectors
are strictly non-negative, the iteration avoids catastrophic subtractive
cancellation. A standard forward-error estimate suggests an accumulated
relative rounding error on the order of
$n_{\max} \cdot m \cdot \varepsilon \approx 6.7 \times 10^{-13}$ (at
$m=15, n=200$, with machine epsilon $\varepsilon \approx 2.22 \times 10^{-16}$).
Comparisons use a relative tolerance of $10^{-11}$, which safely exceeds this
estimate by more than an order of magnitude. It is not an outward-rounded
error enclosure. This floating-point engine acts strictly as a high-speed
filter: all flagged anomalies are subsequently verified using exact
integer arithmetic. The filter can miss a true crossover whose relative margin
lies below the tolerance, so its negative results are conditional screens, not
unconditional nonexistence theorems.
The total cost per target graph is $O(n_{\max} \cdot m^2)$, giving
$O(h \cdot n_{\max} \cdot m^2)$ for the full sweep over $h$ targets.

Throughput: $\sim\!256{,}000$ graphs/second on six threads (Intel
Coffee Lake, AVX2 FMA) using dense $16 \times 16$ padded adjacency
matrices via optimized DAXPY loop
ordering and OpenMP parallelism. The bipartite sweep
of 608,930,559 connected bipartite graphs on $2 \le m \le 15$ vertices
(validated against OEIS A005142~\cite{oeis}) completes in approximately
35~minutes.

\begin{remark}[Depth-dependent bipartite Leontovich graph]
\label{rem:depth_anomaly}
While no $d = 2$ Leontovich violations were found, the $m = 15$
bipartite sweep revealed a single graph $H^*$ (Figure~\ref{fig:anomaly})
for which $\Hom(E_n^{(d)}, H^*) < \Hom(P_n, H^*)$ at several larger
depths, verified with exact integer arithmetic.
This bipartite graph on 15 vertices with partition sizes $(3,12)$
and degree sequence $[7, 5, 5, 2^5, 1^7]$ exhibits 300
such anomalies across $d \in \{14,16,18,20\}$ and odd
$n \le 199$; the first occurs at $(n,d)=(49,16)$ with
\[
  \Hom(P_{49},H^*)-\Hom(E_{49}^{(16)},H^*) =
  41{,}377{,}493{,}496{,}440{,}451{,}164 > 0.
\]
The maximum relative deficit in the tested range is
$\sim 7.6 \times 10^{-4}$.
No anomalies of any depth were \emph{flagged} at $m \le 14$ in this
floating-point screen; as above, this is conditional on the tolerance and does
not exclude an unflagged crossover with a smaller relative margin.
\end{remark}

\begin{figure}[H]
  \centering
  \includegraphics[width=0.38\textwidth]{anomaly_NCQCCA.png}%
  \hfill
  \includegraphics[width=0.58\textwidth]{anomaly_ratio.png}
  \caption{\textbf{Left:} The depth-anomaly graph $H^*$ on 15
    vertices\protect\footnotemark; red/blue shows the $(3,12)$ bipartition.
    \textbf{Right:} Ratio $\Hom(E_n^{(d)}, H^*) / \Hom(P_n, H^*)$ for
    odd~$n$. For the displayed larger depths the ratio drops below~1
    (shaded), meaning $E_n^{(d)}$ beats the path. For $d = 2$ the path
    remains the minimizer.}
  \label{fig:anomaly}
\end{figure}
\footnotetext{Graph6 encoding: \texttt{NCQCCA?\_B?K?W?g?K??}}

To search beyond the exhaustive enumeration wall,
we employ a simulated annealing search seeded from $T(7,1,9)$,
mutating edges and contracting vertices to seek smaller Leontovich
graphs within the graph manifold.

\begin{remark}
Since an independent set in $G$ is exactly a homomorphism
$G \to H_{\mathrm{ind}}$ (two vertices, exactly one of which has a loop, connected by an edge), the
framework developed here extends naturally to the Merrifield--Simmons
index: degree-constrained trees that maximize independent sets are
dual to trees that maximize homomorphisms into the same looped target.
\end{remark}
